\documentclass[../main.tex]{subfiles}
\begin{document}

\chapter{Cache Coherence}
\lessoninfo{
  video={Lesson 19, videos 1--33},
  textbook={H\&P \cite{hennessy2017} §5.2--5.4},
  keywords={cache coherence, bus snooping, write-update coherence, write-invalidate coherence, MSI protocol, cache-to-cache transfer, MOSI protocol, MESI protocol, MOESI protocol, directory-based coherence, coherence miss, true sharing, false sharing},
}

In the shared-memory multiprocessors of Lesson 18 every core has caches of
its own, while all cores share one main memory.  A cache holds a private copy
of the data it contains: when one core writes to its copy, the copies in the
other caches---and, under the write-back policy, main memory---still hold the
old value.  This lesson shows how such stale copies make a multiprocessor
return wrong results, defines the property that rules them out, cache
coherence, and develops the mechanisms that provide it: write-update and
write-invalidate coherence on a shared bus, the MSI protocol and its
refinements MOSI, MESI and MOESI, and directory-based coherence for larger
machines.  It ends with the additional cache misses that coherence causes.

\section{The coherence problem}
\label{sec:l19-problem}

Consider a multiprocessor with centralized shared memory
(\cref{sec:l18-centralized}) in which each core has a private cache that
follows the write-back policy of Lesson 12.\source{Lesson 19, videos 1--3;
H\&P §5.2.}  In a uniprocessor this policy is harmless: main memory may hold
an out-of-date value for a while, but the only processor that could read it
finds the newer copy in its own cache first.  With several caches this is no
longer true.  \Cref{fig:l19-problem} follows a memory location X that
initially holds 5.
\begin{enumerate}
  \item Core 0 reads X.  The read misses, and core 0's cache loads the block
    that contains X from memory.
  \item Core 1 reads X and likewise loads a copy of the block.
  \item Core 0 writes 8 to X.  The write hits and, under write-back, changes
    only the copy in core 0's cache, which becomes dirty.
  \item Core 1 reads X.  The read hits in core 1's cache and returns 5,
    although 8 has been written.
  \item Core 2 reads X.  The read misses and is served by main memory, which
    also still holds 5.
\end{enumerate}

\begin{figure}[htbp]
  \centering
  % The cache coherence problem: three cores with private write-back caches
% share one main memory.  Core 0 has overwritten X in its cache only; core 1
% reads a stale copy from its cache, core 2 a stale value from memory.
\begin{tikzpicture}[x=1cm,y=1cm, font=\footnotesize\sffamily, >={Stealth[length=1.8mm]},
  core/.style={draw, fill=ln-accent!15, minimum width=2.5cm, minimum height=0.6cm, inner sep=2pt},
  cache/.style={draw, fill=ln-box, minimum width=2.5cm, minimum height=0.6cm, inner sep=2pt},
  big/.style={draw, fill=ln-box, minimum height=0.7cm, inner sep=3pt},
  op/.style={font=\scriptsize, align=center, anchor=south},
  bad/.style={text=ln-caution}]
  \foreach \i/\c in {0/{X $=$ 8\enspace\iflangja{（ダーティ）}{(dirty)}},
                     1/{X $=$ 5},
                     2/{\iflangja{X のコピーなし}{no copy of X}}} {
    \node[core] (c\i) at (\i*3.4,0) {\iflangja{コア}{core}~\i};
    \node[cache] (k\i) at (\i*3.4,-0.85) {\c};
    \draw (c\i.south) -- (k\i.north);
    \draw (k\i.south) -- (\i*3.4,-1.75);
  }
  % operations, numbered in the order they happen
  \node[op] at (c0.north) {\iflangja{(1) X を読む\\(3) X に 8 を書く}%
                                   {(1) read X\\(3) write 8 to X}};
  \node[op] at (c1.north) {\iflangja{(2) X を読む}{(2) read X}\\
    \textcolor{ln-caution}{\iflangja{(4) X を読む: ヒット，5}%
                                   {(4) read X: hit, 5}}};
  \node[op, bad] at (c2.north) {\iflangja{(5) X を読む:\\ミス，主記憶から 5}%
                                        {(5) read X:\\miss, 5 from memory}};
  % bus and memory
  \draw[line width=1.4pt] (-1.45,-1.75) -- (8.25,-1.75)
    node[right, font=\footnotesize] {\iflangja{バス}{bus}};
  \node[big, minimum width=4.2cm] (mem) at (3.4,-2.75) {\iflangja{主記憶: X $=$ 5}{main memory: X $=$ 5}};
  \draw (mem.north) -- (3.4,-1.75);
\end{tikzpicture}

  \caption{Incoherent caches.  Cores 0 and 1 read X, which holds 5 (steps 1
    and 2), and core 0 writes 8 into its cache (step 3).  The later reads of
    core 1 (step 4) and core 2 (step 5) both return the stale value 5.}
  \label{fig:l19-problem}
\end{figure}

The two wrong reads have different causes.  Step 5 reads a stale value from
memory; the write-through policy, which updates memory on every write, would
avoid it.  Step 4 reads a stale copy in another cache, and no write policy of
an individual cache can avoid that: however core 0's cache handles the write,
the copy in core 1's cache does not change.\caution{Write-through does not
make private caches coherent: the copies in the other caches still go
stale.}

\section{Coherence}
\label{sec:l19-definition}

A shared-memory program expects the behavior of a single memory that all
cores share, without caches.\source{Lesson 19, videos 4--5; H\&P §5.2.}

\begin{definition}[Cache coherence]
  A multiprocessor whose cores have private caches provides
  \term{cache coherence}[キャッシュコヒーレンス] if the following hold for
  every memory location X.
  \begin{enumerate}
    \item A read of X on a core returns the value most recently written to X
      by that core, if no other core has written X in between.
    \item If one core writes X and another core reads X sufficiently later,
      with no write to X in between, the read returns the value written by
      the first core.
    \item Writes to X are serialized: all cores observe any two writes to X
      in the same order.
  \end{enumerate}
\end{definition}

The first condition is the behavior of memory in a uniprocessor, which every
cache already preserves.  The second extends it to other cores but allows a
delay: a read issued while another core's write is still in progress may
return the old value.  The third lets all cores agree which of two nearly
simultaneous writes came last.  The reads of steps 4 and 5 in
\cref{fig:l19-problem} violate the second condition: the caches are
\emph{incoherent}.

\begin{note}
  Coherence is a requirement on what reads return, not on what the copies
  contain.  Copies may disagree without violating it: main memory may hold an
  older value than a dirty block in a cache, and a cache may still contain an
  old value in a line that is marked invalid.  The system is coherent as long
  as every read is served from a copy that holds the latest value.
\end{note}

\section{Achieving coherence}
\label{sec:l19-achieving}

The definition and the example of \cref{fig:l19-problem} show what a
mechanism for coherence has to do.\source{Lesson 19, video 6; H\&P §5.2.}
\begin{enumerate}
  \item \textbf{Make writes visible.}  Under write-back a write that hits
    causes no activity outside the core, so the other caches cannot know that
    their copies have become stale.  Writes to blocks that other caches may
    hold must therefore be announced to them.
  \item \textbf{Act on the writes of other cores.}  A cache that holds a
    copy of a block written by another core must either \emph{update} its copy
    with the new value or \emph{invalidate} it, so that its next access to the
    block misses and fetches the new value.
  \item \textbf{Order the writes.}  Writes to the same block by different
    cores must be applied to all copies in one order that all caches
    observe.
  \item \textbf{Serve reads from the latest copy.}  A read that misses must
    obtain the latest value, which may be in another cache rather than in
    memory.
\end{enumerate}

With centralized shared memory, all caches and main memory are connected to
one bus, and every operation on the bus reaches every device connected to
it.  Coherence can therefore be built on
\term{bus snooping}[バススヌーピング]: every cache monitors all operations on
the bus, including those addressed to memory or to other caches, and reacts
to those that concern blocks it holds.  Snooping also orders the operations,
because the bus carries one operation at a time and every cache observes them
in that order.  Machines without a shared bus need another mechanism for
both tasks, the directory of \cref{sec:l19-directory}.  A coherence mechanism
thus combines two choices: other copies are either updated or invalidated,
and writes are observed and ordered either by snooping or by a directory.

\section{Write-update coherence}
\label{sec:l19-write-update}

The first way of acting on a write is to propagate the new value to every
copy.\source{Lesson 19, videos 7--11; H\&P §5.2.}

\begin{definition}[Write-update coherence]
  Under \term{write-update coherence}[書き込み時更新方式] every write is
  broadcast with its address and its new value, and every other cache that
  holds a copy of the written block updates its copy.
\end{definition}

With snooping, each core simply places every write on the bus.  Main memory
performs the write like any other, and the other caches snoop it and update
their copies if they hold the block.  In \cref{fig:l19-problem} the write of
step 3 would then change X to 8 in memory and in core 1's cache, and both
later reads would return 8.  Because the bus carries one operation at a time,
writes by different cores are applied to all copies in the same order, even
when two cores write the same location at the same moment, as the third
condition of coherence requires.

This basic scheme is correct but expensive: every write occupies the shared
bus and is also performed in main memory---exactly the cost that the
write-back policy is meant to avoid.  Two optimizations remove most of it.

\subsection{Avoiding memory writes}

The dirty bit of Lesson 12 lets a cache defer the memory write.  A write is
still placed on the bus so that the other caches can update their copies,
but memory ignores it; instead, the writing cache sets the dirty bit of the
block and writes the block back to memory when it evicts it.  A cache that
snoops another core's write to a block it holds dirty updates its copy and
clears its dirty bit, so only the last writer is responsible for the block.
While a dirty copy exists, memory is stale.  The cache with the dirty copy
therefore also snoops the bus for \emph{reads} of that block: when another
core misses on the block and places the read on the bus, this cache supplies
the data, and memory does not respond.

\subsection{Avoiding bus writes}

Most blocks are used by only one core, and broadcasting the writes to such
blocks occupies the bus for nothing.  A \emph{shared bit} in each cache line
records whether other caches may hold the block.  When a cache snoops a read
of a block that it holds, it sets its shared bit and signals on a
\emph{shared line} of the bus that a copy exists; the reading cache sets its
own shared bit if this signal is present.  A write to a block whose shared bit
is clear is performed locally, as in a uniprocessor; only writes to blocks
whose shared bit is set are placed on the bus.\margin{The shared bit is
conservative: it stays set after the other copies have been evicted.}

\begin{example}
  \label{ex:l19-write-update}
  Two cores access a block X that no cache holds initially.
  \Cref{tab:l19-write-update} counts the bus operations and memory writes of
  six accesses under the basic scheme and with the optimizations added.  With
  the dirty bit, memory is not written during the sequence (core 0 writes the
  block back once, when it evicts it), and the read of step 4 is served by
  core 0.  With the shared bit in addition, the two writes that occur before
  core 1 has a copy stay local, and only the write of step 5 is broadcast.
\end{example}

\begin{table}[htbp]
  \centering\small
  \caption{Write-update coherence (\cref{ex:l19-write-update}).  rd: read on
    the bus; wr: write on the bus or into memory; rd$^{*}$: read served by
    core 0 instead of memory.}
  \label{tab:l19-write-update}
  \begin{tabular}{@{}cl cc cc cc@{}}
    \toprule
    & & \multicolumn{2}{c}{Basic} & \multicolumn{2}{c}{Dirty bit}
      & \multicolumn{2}{c}{Dirty and shared bits} \\
    \cmidrule(lr){3-4}\cmidrule(lr){5-6}\cmidrule(l){7-8}
    Step & Access & Bus & Memory & Bus & Memory & Bus & Memory \\
    \midrule
    1 & core 0 reads X  & rd  & --- & rd       & --- & rd       & --- \\
    2 & core 0 writes X & wr  & wr  & wr       & --- & ---      & --- \\
    3 & core 0 writes X & wr  & wr  & wr       & --- & ---      & --- \\
    4 & core 1 reads X  & rd  & --- & rd$^{*}$ & --- & rd$^{*}$ & --- \\
    5 & core 0 writes X & wr  & wr  & wr       & --- & wr       & --- \\
    6 & core 1 reads X  & --- & --- & ---      & --- & ---      & --- \\
    \midrule
    \multicolumn{2}{@{}l}{Total} & 5 & 3 & 5 & 0 & 3 & 0 \\
    \bottomrule
  \end{tabular}
\end{table}

\section{Write-invalidate coherence}
\label{sec:l19-write-invalidate}

The second way of acting on a write is to remove the other copies instead of
updating them.\source{Lesson 19, videos 12--15; H\&P §5.2.}

\begin{definition}[Write-invalidate coherence]
  Under \term{write-invalidate coherence}[書き込み時無効化方式] a core that
  writes to a block which other caches may hold first broadcasts an
  invalidation of the block, and every other cache that holds a copy of the
  block discards it by clearing its valid bit.
\end{definition}

An invalidation carries only the address of the block.  Once it has been
broadcast, the writing cache holds the only valid copy, so further writes to
the block---to any address within it---are local and need no bus operation
until another core accesses the block.  That access misses, because its copy
has been invalidated, and obtains the latest value either from memory, if
the block has been written back, or from the cache that wrote it.

\subsection{Update versus invalidate}

Which scheme needs fewer bus operations depends on the pattern of accesses.
\begin{itemize}
  \item \textbf{Several writes to a block before another core reads it.}
    Write-update broadcasts every write; write-invalidate broadcasts a single
    invalidation.  This holds for repeated writes to one address and equally
    for writes to different addresses in the same block, since invalidation
    works per block.
  \item \textbf{A thread moves to another core.}  The blocks that the thread
    used remain in the cache of its old core, and when the thread reads them
    on its new core, both caches mark them shared.  Under write-update, every
    later write of the thread to such a block is broadcast and updates a copy
    that the old core no longer uses, until that copy happens to be evicted.
    Under write-invalidate, the first write removes the old copy, and all
    further writes are local.
  \item \textbf{Writes and reads by different cores alternate.}  If one core
    writes a block and another core reads the new value after every write,
    write-update needs one bus write per write, and the reads all hit.
    Write-invalidate needs an invalidation per write and a miss per read.
\end{itemize}
Only the last pattern favors write-update, and the first two are the common
case.  Multiprocessors therefore use write-invalidate coherence, and all
protocols in the rest of this lesson are based on it.

\begin{example}
  Cores 0 and 1 both hold a copy of block X.  If core 0 writes X ten times
  and core 1 then reads X once, write-update needs 10 bus writes, and core
  1's read hits;
  write-invalidate needs one invalidation and one bus read for core 1's miss,
  2 bus operations in total.  If instead core 1 reads X after each of the ten
  writes, write-update still needs 10 bus operations, whereas write-invalidate
  needs 10 invalidations and 10 misses, 20 in total.
\end{example}

\begin{keypoint}
  Write-update keeps every copy current at the price of a bus operation per
  write; write-invalidate pays one invalidation per block and then lets the
  writing core proceed locally.  Because writes usually come in bursts from
  one core, practical coherence protocols invalidate.
\end{keypoint}

\section{The MSI protocol}
\label{sec:l19-msi}

Write-invalidate coherence with snooping is described by a small state
machine for each cache line.\source{Lesson 19, videos 16--19; H\&P §5.2.}
The state of the block in a line changes with the accesses of the core itself
and with the accesses of other cores that the cache observes on the bus.  The
simplest such state machine, the \term{MSI protocol}[MSI プロトコル], uses
the valid and dirty bits of a line to distinguish three states.
\begin{description}
  \item[M (modified)] Valid and dirty.  The block has been written in this
    cache, and all other copies have been invalidated: this cache holds the
    only valid copy, which is newer than the block in memory.  Reads and
    writes are local.
  \item[S (shared)] Valid and not dirty.  Other caches may hold valid copies
    as well.  Reads are local; a write must first broadcast an invalidation.
  \item[I (invalid)] The valid bit is clear.  Every access misses.
\end{description}

\begin{figure}[htbp]
  \centering
  % MSI state machine of one block in one cache.  Solid edges: access by this
% core; dashed edges: access by another core, observed on the bus.
\begin{tikzpicture}[x=1cm,y=1cm, font=\footnotesize\sffamily, >={Stealth[length=1.8mm]},
  st/.style={circle, draw, minimum size=0.9cm, inner sep=1pt, fill=ln-box},
  dirty/.style={st, fill=ln-accent!15},
  lab/.style={font=\scriptsize, inner sep=2pt},
  other/.style={dashed, ln-accent}]
  \node[st]    (I) at (0,0)     {I};
  \node[st]    (S) at (3.4,-1.8) {S};
  \node[dirty] (M) at (6.8,0)   {M};
  % this core
  \draw[->] (I) to[bend right=22] node[lab, below left]
    {\iflangja{R / バス読み出し}{R / bus read}} (S);
  \draw[->] (I) -- node[lab, above]
    {\iflangja{W / バス読み出し＋無効化}{W / bus read + invalidate}} (M);
  \draw[->] (S) to[bend left=22] node[lab, above left, pos=0.3]
    {\iflangja{W / 無効化}{W / invalidate}} (M);
  % other cores
  \draw[->, other] (S) to[bend right=22] node[lab, above right, pos=0.45] {W$'$} (I);
  \draw[->, other] (M) to[bend left=22] node[lab, below right]
    {\iflangja{R$'$ / 書き戻し}{R$'$ / write back}} (S);
  \draw[->, other] (M) to[bend right=30] node[lab, above]
    {\iflangja{W$'$ / 書き戻し}{W$'$ / write back}} (I);
\end{tikzpicture}

  \caption{The MSI protocol for one block in one cache.  R, W: read, write by
    this core; R$'$, W$'$ (dashed): read, write by another core, observed on
    the bus.  Accesses without an edge leave the state unchanged and need no
    bus operation.  The shaded state holds a block newer than memory.}
  \label{fig:l19-msi}
\end{figure}

\Cref{fig:l19-msi} shows the transitions.  Accesses by the core itself
(solid edges) are handled as follows.
\begin{itemize}
  \item A read in I misses: the cache places the read on the bus, loads the
    block and enters S.
  \item A write in I misses too: in one bus operation, the cache reads the
    block and invalidates all other copies, and it enters M.
  \item A write in S broadcasts an invalidation and enters M.
  \item Reads in S, and reads and writes in M, need no bus operation.
\end{itemize}
Accesses by other cores (dashed edges) are observed on the bus.
\begin{itemize}
  \item A write by another core, announced by its invalidation, takes S to
    I.  The copy is discarded without being written back, since memory holds
    the same value.
  \item In M the cache holds the only up-to-date value of the block.  When
    another core reads the block, the cache writes the block back to memory
    and enters S, because the other core now holds a copy too.  When another
    core writes the block, the cache writes the block back to memory, from
    which the writer obtains it, and enters I.
\end{itemize}
When a cache evicts a block, it writes the block back to memory only in M; a
block in S is simply discarded.

\subsection{Cache-to-cache transfers}
\label{sec:l19-c2c}

When a core misses on a block that another cache holds in M, whether on a
read or on a write, memory would answer the bus read with a stale value.  The
cache in M, which snoops the read, can prevent this in two ways.
\begin{enumerate}
  \item \emph{Abort and retry}: it aborts the read and writes the block back
    to memory; the reading core then issues the read again, and memory serves
    it with the current value.  This costs a memory write and a memory read.
  \item \emph{Intervention}: it signals that memory must not answer and
    places the block on the bus itself.  On a read miss, both copies end in
    S, which is not dirty, so memory takes the block from the bus as well.
\end{enumerate}
Intervention, a \term{cache-to-cache transfer}[キャッシュ間転送], needs
more complex bus hardware, but a single transfer serves the miss and, on a
read miss, also performs the write-back that MSI requires on the transition
from M to S.  The dirty-bit optimization of write-update coherence
(\cref{sec:l19-write-update}) answers reads in the same way, but there the
cache stays dirty and memory does not take the block.

\begin{example}
  \label{ex:l19-msi}
  Three cores access a block X that no cache holds initially
  (\cref{tab:l19-msi}).  The reads of steps 1 and 2 leave two copies in S.
  Core 1's write in step 3 invalidates core 0's copy, and its second write
  (step 4) is local.  When core 2's read appears on the bus in step 5, core 1
  holds the only current value; it writes the block back to memory, and both
  copies end in S.  Core 0's write in step 6 misses and invalidates both of
  them.  Five of the seven accesses need a bus operation, and one causes a
  memory write.
\end{example}

\begin{table}[htbp]
  \centering\small
  \caption{Block X under the MSI protocol (\cref{ex:l19-msi}): state of the
    block in each cache after the access.}
  \label{tab:l19-msi}
  \begin{tabular}{@{}cl ccc l@{}}
    \toprule
    Step & Access & Core 0 & Core 1 & Core 2 & Bus and memory \\
    \midrule
    1 & core 0 reads X  & S & I & I & read, served by memory \\
    2 & core 1 reads X  & S & S & I & read, served by memory \\
    3 & core 1 writes X & I & M & I & invalidation \\
    4 & core 1 writes X & I & M & I & --- \\
    5 & core 2 reads X  & I & S & S & read; core 1 writes X back \\
    6 & core 0 writes X & M & I & I & read and invalidation \\
    7 & core 0 reads X  & M & I & I & --- \\
    \bottomrule
  \end{tabular}
\end{table}

\section{The MOSI protocol}
\label{sec:l19-mosi}

A cache-to-cache transfer does not remove the memory write of the transition
from M to S.\source{Lesson 19, videos 20--21.}  After the transfer both
copies are in S, which is not dirty, so both caches would discard their
copies on eviction; unless memory takes the block during the transfer, the
latest value is lost once both copies are gone.  A block that one core writes
and another core reads in alternation therefore costs a memory write for
every write, although neither cache evicts it.

The \term{MOSI protocol}[MOSI プロトコル] avoids these writes by making one
of the copies of a shared block its owner.
\begin{description}
  \item[O (owned)] Valid; other caches may hold copies in S.  The block is
    newer than the block in memory, and this cache is responsible for it: it
    answers reads of the block by other cores and writes the block back to
    memory when it evicts it.  Reads are local; a write broadcasts an
    invalidation and enters M.
\end{description}

\begin{figure}[htbp]
  \centering
  % MOSI state machine: MSI plus the owned state O.  Solid edges: access by
% this core; dashed edges: access by another core, observed on the bus.
\begin{tikzpicture}[x=1cm,y=1cm, font=\footnotesize\sffamily, >={Stealth[length=1.8mm]},
  st/.style={circle, draw, minimum size=0.9cm, inner sep=1pt, fill=ln-box},
  dirty/.style={st, fill=ln-accent!15},
  lab/.style={font=\scriptsize, inner sep=2pt},
  other/.style={dashed, ln-accent}]
  \node[st]    (I) at (0,0)      {I};
  \node[st]    (S) at (3.4,-1.8) {S};
  \node[dirty] (M) at (6.8,0)    {M};
  \node[dirty] (O) at (9.8,0)    {O};
  % this core
  \draw[->] (I) to[bend right=22] node[lab, below left]
    {\iflangja{R / バス読み出し}{R / bus read}} (S);
  \draw[->] (I) -- node[lab, above]
    {\iflangja{W / バス読み出し＋無効化}{W / bus read + invalidate}} (M);
  \draw[->] (S) to[bend left=22] node[lab, above left, pos=0.3]
    {\iflangja{W / 無効化}{W / invalidate}} (M);
  \draw[->] (O) to[bend left=30] node[lab, below]
    {\iflangja{W / 無効化}{W / invalidate}} (M);
  % other cores
  \draw[->, other] (S) to[bend right=22] node[lab, above right, pos=0.45] {W$'$} (I);
  \draw[->, other] (M) to[bend right=30] node[lab, above]
    {\iflangja{W$'$ / ブロックを供給}{W$'$ / supply block}} (I);
  \draw[->, other] (M) to[bend left=30] node[lab, above]
    {\iflangja{R$'$ / ブロックを供給}{R$'$ / supply block}} (O);
  \draw[->, other] (O) to[out=-60, in=-120, looseness=5] node[lab, below]
    {\iflangja{R$'$ / ブロックを供給}{R$'$ / supply block}} (O);
  \draw[->, other, rounded corners=6pt] (O.north) -- ++(0,1.75) -| node[lab, above, pos=0.25]
    {W$'$} (I.north);
\end{tikzpicture}

  \caption{The MOSI protocol.  Notation as in \cref{fig:l19-msi}; the shaded
    states are responsible for a block newer than memory.  On W$'$ in O, the
    block is supplied only if the writer misses.}
  \label{fig:l19-mosi}
\end{figure}

When another core reads a block that the cache holds in M, the cache supplies
the block by a cache-to-cache transfer and enters O instead of S, without
writing memory (\cref{fig:l19-mosi}).  In O it supplies the block to every
further reader.  A write by the owner broadcasts an invalidation and returns
to M; a write by another core takes O to I, and only if the writer missed
does the owner first supply the block, since a writer in S already holds the
current value.  The block reaches memory only when a cache evicts it in
M or O.  In \cref{ex:l19-msi}, MOSI would take core 1 to O instead of S in
step 5, without the memory write, and core 1 would supply the block for core
0's write in step 6.

\section{The exclusive state}
\label{sec:l19-mesi}

MSI and MOSI handle a very common access pattern
inefficiently.\source{Lesson 19, videos 22--25; H\&P §5.2;
\cite{papamarcos1984}.}  Programs often read a variable and then write it, as
in an increment, and most variables are used by only one thread.  For such a
block the read misses and enters S, and the write then broadcasts an
invalidation---although no other cache holds a copy.  Every such read-write
pair needs two bus operations, where a uniprocessor has only the miss.

S covers two cases that the cache cannot tell apart: other caches hold
copies, or the block is unshared but not yet written.  A fourth state
separates them.
\begin{description}
  \item[E (exclusive)] Valid and not dirty; no other cache holds a copy.
    Reads are local, and so is a write, which enters M without any bus
    operation.
\end{description}
A read in I enters E if no other cache holds the block, and S otherwise.  The
cache learns which case applies during the bus read itself: caches that hold
the block snoop the read and signal on the shared line of the bus, as in
\cref{sec:l19-write-update}.  When another core reads a block held in E, the
cache enters S without a write, since memory is current.  A write by another
core takes E to I.

Adding E to MSI gives the \term{MESI protocol}[MESI プロトコル], and adding
it to MOSI gives the \term{MOESI protocol}[MOESI プロトコル]
(\cref{fig:l19-moesi}).  \Cref{tab:l19-states} summarizes the five states.

\begin{figure}[htbp]
  \centering
  % MOESI state machine: MOSI plus the exclusive state E.  Solid edges: access
% by this core; dashed edges: access by another core, observed on the bus.
% A write by another core takes every valid state to I (edges not drawn).
\begin{tikzpicture}[x=1cm,y=1cm, font=\footnotesize\sffamily, >={Stealth[length=1.8mm]},
  st/.style={circle, draw, minimum size=0.9cm, inner sep=1pt, fill=ln-box},
  dirty/.style={st, fill=ln-accent!15},
  lab/.style={font=\scriptsize, inner sep=2pt},
  other/.style={dashed, ln-accent}]
  \node[st]    (I) at (0,0)      {I};
  \node[st]    (E) at (3.4,1.8)  {E};
  \node[st]    (S) at (3.4,-1.8) {S};
  \node[dirty] (M) at (6.8,0)    {M};
  \node[dirty] (O) at (9.8,0)    {O};
  % this core
  \draw[->] (I) -- node[lab, above, sloped]
    {\iflangja{R，他にコピーなし}{R, no other copy}} (E);
  \draw[->] (I) -- node[lab, below, sloped]
    {\iflangja{R，他にコピーあり}{R, other copies}} (S);
  \draw[->] (I) -- node[lab, above, pos=0.74, align=center]
    {\iflangja{W / バス読み出し\\＋無効化}{W / bus read\\+ invalidate}} (M);
  \draw[->] (E) -- node[lab, above, sloped] {W} (M);
  \draw[->] (S) -- node[lab, below, sloped]
    {\iflangja{W / 無効化}{W / invalidate}} (M);
  \draw[->] (O) to[bend left=30] node[lab, below]
    {\iflangja{W / 無効化}{W / invalidate}} (M);
  % other cores
  \draw[->, other] (E) -- node[lab, right, pos=0.25] {R$'$} (S);
  \draw[->, other] (M) to[bend left=30] node[lab, above]
    {\iflangja{R$'$ / ブロックを供給}{R$'$ / supply block}} (O);
  \draw[->, other] (O) to[out=-60, in=-120, looseness=5] node[lab, below]
    {\iflangja{R$'$ / ブロックを供給}{R$'$ / supply block}} (O);
\end{tikzpicture}

  \caption{The MOESI protocol.  Notation as in \cref{fig:l19-msi}.  Every
    transition out of I performs a bus read.  A write by another core takes
    every valid state to I; if the writer misses, a cache in M or O supplies
    the block (these edges are not drawn).}
  \label{fig:l19-moesi}
\end{figure}

\begin{table}[htbp]
  \centering\small
  \caption{The states of MOESI.  MSI uses M, S and I; MOSI adds O and MESI
    adds E.  A responsible cache answers reads by other cores and writes the
    block back on eviction.}
  \label{tab:l19-states}
  \begin{tabular}{@{}lcccc@{}}
    \toprule
    State & Valid & Copies elsewhere & Responsible & Write without bus \\
    \midrule
    M (modified)  & yes & no       & yes & yes \\
    O (owned)     & yes & possible & yes & no \\
    E (exclusive) & yes & no       & no  & yes \\
    S (shared)    & yes & possible & no  & no \\
    I (invalid)   & no  & possible & no  & no \\
    \bottomrule
  \end{tabular}
\end{table}

\begin{example}
  \label{ex:l19-protocols}
  Two cores access a block X that no cache holds initially: core 0 reads and
  then writes X, core 1 reads it, core 0 writes it again, and core 1 reads it
  again.  \Cref{tab:l19-protocols} gives the states and operations under the
  four protocols.  E saves the invalidation of step 2, because core 0's read
  found no other copy; O saves the memory writes of steps 3 and 5, at the
  price of one write when core 0 eventually evicts X.
\end{example}

\begin{table}[htbp]
  \centering\small
  \caption{Block X under four protocols (\cref{ex:l19-protocols}): states in
    core 0 and core 1 after each access.  b: bus operation; m: memory write.}
  \label{tab:l19-protocols}
  \begin{tabular}{@{}cl cc cc cc cc@{}}
    \toprule
    & & \multicolumn{2}{c}{MSI} & \multicolumn{2}{c}{MOSI}
      & \multicolumn{2}{c}{MESI} & \multicolumn{2}{c}{MOESI} \\
    \cmidrule(lr){3-4}\cmidrule(lr){5-6}\cmidrule(lr){7-8}\cmidrule(l){9-10}
    Step & Access & States & Ops & States & Ops & States & Ops & States & Ops \\
    \midrule
    1 & core 0 reads X  & S\,I & b    & S\,I & b   & E\,I & b    & E\,I & b \\
    2 & core 0 writes X & M\,I & b    & M\,I & b   & M\,I & ---  & M\,I & --- \\
    3 & core 1 reads X  & S\,S & b, m & O\,S & b   & S\,S & b, m & O\,S & b \\
    4 & core 0 writes X & M\,I & b    & M\,I & b   & M\,I & b    & M\,I & b \\
    5 & core 1 reads X  & S\,S & b, m & O\,S & b   & S\,S & b, m & O\,S & b \\
    \midrule
    \multicolumn{2}{@{}l}{Bus operations} & & 5 & & 5 & & 4 & & 4 \\
    \multicolumn{2}{@{}l}{Memory writes}  & & 2 & & 0 & & 2 & & 0 \\
    \bottomrule
  \end{tabular}
\end{table}

\begin{keypoint}
  In all four protocols only a cache in M or O is responsible for a block
  newer than memory, and only a cache in M or E may write without a bus
  operation.  The two additions are independent: the E state saves the
  invalidation after reading unshared data, and the O state saves the memory
  write when a modified block is read by another core.
\end{keypoint}

\section{Directory-based coherence}
\label{sec:l19-directory}

Snooping requires every cache to observe every bus operation, which ties it
to a shared bus.\source{Lesson 19, videos 26--30; H\&P §5.4;
\cite{agarwal1988}.}  A shared bus limits a machine to a few cores (Lesson
18), and a larger machine whose cores share one address space over a network
(distributed shared memory in the sense of H\&P, see the terminology note in
Lesson 18) has no bus at all.  Sending every read and invalidation to all
caches over the network would cost far more than the operations themselves,
and without a bus nothing orders the operations of different cores on a
block.  Larger machines therefore keep explicit track of where each block is
cached.

\begin{definition}[Directory-based coherence]
  Under \term{directory-based coherence}[ディレクトリ方式] a
  \emph{directory} records, for every block, which caches hold a copy of it.
  A cache sends each operation that a snooping protocol would broadcast to the
  directory, which forwards it only to the caches that hold the block.  All
  such requests for a block pass through its directory entry, which thereby
  serializes them.
\end{definition}

A directory entry needs little information (\cref{fig:l19-directory}):
\begin{itemize}
  \item an \emph{exclusive} bit, which indicates that a single cache holds
    the block, in M or E, and can write it without notifying anyone;
  \item one \emph{presence} bit per cache, which indicates whether that cache
    holds a valid copy.
\end{itemize}
The states themselves are still kept in the caches, so the protocols of the
previous sections can be used unchanged.  Accesses that a protocol handles
locally---reads in every valid state, writes in M and E---do not involve the
directory; the others are sent to it.
\begin{itemize}
  \item \textbf{Read miss.}  If no cache holds the block, memory supplies
    it, and the reader becomes its only holder (in E under MESI and MOESI,
    with the exclusive bit set).  Otherwise the directory forwards the read to
    a cache that holds the block, which supplies it by a cache-to-cache
    transfer; if the entry is exclusive, the read goes to the exclusive
    holder, which changes its state as on a snooped read, and the exclusive
    bit is cleared.  In both cases the reader is added to the entry.
  \item \textbf{Write in S or O, or write miss.}  The directory forwards the
    invalidation only to the caches whose presence bits are set and records
    the writer as the exclusive holder.  On a write miss, a cache that holds
    the block supplies it before entering I, or memory does if no cache holds
    it.
\end{itemize}
When a cache evicts a block, writing it back in M or O, it notifies the
directory, which clears the presence bit of that cache and the exclusive
bit, so that the entry remains exact.

The directory need not reside in one place.  It is divided into slices kept
on different nodes, each slice serving a set of blocks, so that no single
node has to handle all requests.  Every request for a block goes to the
slice that serves it, the home of the block, and this slice determines the
order of all accesses to the block.

\begin{figure}[htbp]
  \centering
  % Directory-based coherence: the directory entry of block X records whether
% one cache holds X exclusively and which caches hold a copy.  A write by
% core 3 is sent to the directory, which forwards the invalidation only to
% the caches that hold X.
\begin{tikzpicture}[x=1cm,y=1cm, font=\footnotesize\sffamily, >={Stealth[length=1.8mm]},
  core/.style={draw, fill=ln-accent!15, minimum width=2.1cm, minimum height=0.55cm, inner sep=2pt},
  cache/.style={draw, fill=ln-box, minimum width=2.1cm, minimum height=0.55cm, inner sep=2pt},
  cell/.style={draw, fill=white, minimum width=0.55cm, minimum height=0.5cm, inner sep=0pt},
  hd/.style={font=\scriptsize, text=ln-muted},
  lab/.style={font=\scriptsize, inner sep=2pt},
  msg/.style={->, dashed, ln-accent, thick}]
  \foreach \i/\s in {0/S, 1/O, 2/I, 3/S} {
    \node[core] (c\i) at (\i*2.8,0) {\iflangja{コア}{core}~\i};
    \node[cache] (k\i) at (\i*2.8,-0.75) {\iflangja{キャッシュ}{cache}: X (\s)};
    \draw (c\i.south) -- (k\i.north);
    \draw (k\i.south) -- (\i*2.8,-1.55);
  }
  \draw[line width=1.4pt] (-1.15,-1.55) -- (9.55,-1.55)
    node[right, font=\scriptsize, inner sep=2pt] {\iflangja{ネットワーク}{network}};
  % directory entry of X
  \node[font=\footnotesize] at (2.85,-3.1) {X};
  \foreach \j/\h/\v in {0/Ex/0, 1/C0/1, 2/C1/1, 3/C2/0, 4/C3/1} {
    \node[cell] (e\j) at (3.4+\j*0.55,-3.1) {\v};
    \node[hd, above] at (e\j.north) {\h};
  }
  \fill[ln-accent!15] (e0.south west) rectangle (e0.north east);
  \node[font=\footnotesize] at (e0) {0};
  \draw (e0.south west) rectangle (e0.north east);
  \node[draw=ln-rule, rounded corners=2pt, fit=(e0)(e4), inner xsep=16pt, inner ysep=14pt,
        yshift=2pt, xshift=-5pt] (dir) {};
  \node[lab, text=ln-muted, anchor=north] at (dir.south) {\iflangja{ディレクトリ}{directory}};
  % messages for a write by core 3
  \draw[msg, rounded corners=4pt] (k3.south) ++(0.35,0) |- node[lab, pos=0.25, right]
    {\iflangja{(1) X に書き込み}{(1) write X}} (dir.east);
  \draw[msg, rounded corners=4pt] (dir.west) -| node[lab, pos=0.75, left]
    {\iflangja{(2) 無効化}{(2) invalidate}} ([xshift=0.35cm]k0.south);
  \draw[msg] (dir.north -| k1.south) ++(0.35,0) -- node[lab, pos=0.2, right]
    {\iflangja{(2) 無効化}{(2) invalidate}} ([xshift=0.35cm]k1.south);
\end{tikzpicture}

  \caption{Directory-based coherence in step 5 of \cref{ex:l19-directory}.
    The entry of X shows that cores 0, 1 and 3 hold copies and that none
    holds X exclusively.  Core 3's write (1) is forwarded as an invalidation
    to cores 0 and 1 only (2).}
  \label{fig:l19-directory}
\end{figure}

\needspace{4\baselineskip}
\begin{example}
  \label{ex:l19-directory}
  Four cores use MOESI with a directory (\cref{tab:l19-directory}).  Core 1
  reads and then writes X; its read found no other copy, so X is in E in
  core 1's cache, and the write needs no message.  The read of core 3 goes to
  the directory, which forwards it to the exclusive holder, core 1; core 1
  supplies the block and becomes its owner.  The directory forwards the read
  of core 0 to one of the holders, again core 1.  When core 3 writes X in
  step 5, the directory sends invalidations to cores 0 and 1 only; core 2,
  which holds no copy, receives no message, whereas a broadcast would have
  reached it as well.
\end{example}

\begin{table}[htbp]
  \centering\small
  \caption{Block X with directory-based MOESI (\cref{ex:l19-directory}):
    cache states and directory entry (exclusive bit; presence bits of cores 0
    to 3) after each access, and the messages sent (d: directory).}
  \label{tab:l19-directory}
  \begin{tabular}{@{}cl cccc cc l@{}}
    \toprule
    & & \multicolumn{4}{c}{Cache of core} & \multicolumn{2}{c}{Entry} & \\
    \cmidrule(lr){3-6}\cmidrule(lr){7-8}
    Step & Access & 0 & 1 & 2 & 3 & Ex & Presence & Messages \\
    \midrule
    1 & core 1 reads X  & I & E & I & I & 1 & 0100 & 1$\to$d, memory$\to$1 \\
    2 & core 1 writes X & I & M & I & I & 1 & 0100 & --- \\
    3 & core 3 reads X  & I & O & I & S & 0 & 0101 & 3$\to$d$\to$1$\to$3 \\
    4 & core 0 reads X  & S & O & I & S & 0 & 1101 & 0$\to$d$\to$1$\to$0 \\
    5 & core 3 writes X & I & I & I & M & 1 & 0001 & 3$\to$d$\to$0, 1 \\
    6 & core 3 writes X & I & I & I & M & 1 & 0001 & --- \\
    \bottomrule
  \end{tabular}
\end{table}

\section{Cache misses and coherence}
\label{sec:l19-misses}

Lesson 14 divided cache misses into compulsory, capacity and conflict
misses.  Invalidation adds a fourth kind.\source{Lesson 19, videos 31--33;
H\&P §5.3.}

\begin{definition}[Coherence miss]
  A \term{coherence miss}[コヒーレンスミス] is a miss on a block whose copy in
  the cache has been invalidated because another core wrote the block.
\end{definition}

Coherence misses are of two kinds.
\begin{description}
  \item[True sharing] The core accesses data that the other core actually
    wrote.  Such \term{true sharing}[真共有] makes the miss necessary: without
    the miss, the core would read a stale value.
  \item[False sharing] The other core wrote different data that merely lies
    in the same block.  The data that the core accesses has not changed, and
    the miss is caused only by coherence being maintained per block rather than
    per address.  Such \term{false sharing}[偽共有] increases with the block
    size, since larger blocks contain more unrelated data.
\end{description}

\begin{example}
  \label{ex:l19-false-sharing}
  Threads on cores 0 and 1 each increment a counter of their own, c0 and c1,
  \num{1000} times, and the increments of the two threads alternate.  The
  counters are adjacent and lie in the same cache block
  (\cref{fig:l19-false-sharing}).  Under MOESI, core 0's first increment
  finds no other copy, and its write needs no invalidation.  From core 1's
  first increment on, however, every increment writes the block while the
  other core holds a copy, and invalidates it, so every increment from core
  0's second one on starts with a coherence miss: \num{1998} misses, all due
  to false sharing, because neither thread ever reads the other's counter.
  With the counters in different blocks, only the two compulsory misses
  remain.
\end{example}
\begin{marginfigure}
  \resizebox{\linewidth}{!}{% False sharing: two counters used by different cores lie in one cache block.
\begin{tikzpicture}[x=1cm,y=1cm, font=\footnotesize\sffamily, >={Stealth[length=1.8mm]},
  core/.style={draw, fill=ln-accent!15, minimum width=1.5cm, minimum height=0.55cm, inner sep=2pt},
  cell/.style={draw, minimum width=0.5cm, minimum height=0.5cm, inner sep=0pt}]
  \fill[ln-box] (-0.25,-0.25) rectangle (0.75,0.25);
  \foreach \j in {0,...,7} {
    \node[cell] (b\j) at (\j*0.5,0) {};
  }
  \node[font=\scriptsize] at (b0) {c0};
  \node[font=\scriptsize] at (b1) {c1};
  \node[core] (k0) at (0,1.2)   {\iflangja{コア}{core}~0};
  \node[core] (k1) at (0.5,-1.2) {\iflangja{コア}{core}~1};
  \draw[->] (k0.south) -- (b0.north);
  \draw[->] (k1.north) -- (b1.south);
  \node[font=\scriptsize, text=ln-muted, anchor=north east] at (b7.south east)
    {\iflangja{1つのキャッシュブロック}{one cache block}};
\end{tikzpicture}
}
  \caption{False sharing: counters used by different cores lie in the same
    cache block.}
  \label{fig:l19-false-sharing}
\end{marginfigure}

\begin{keypoint}[Lesson summary]
  The private caches of a shared-memory multiprocessor must be coherent: a
  core sees its own writes, the writes of other cores after a short time,
  and all writes to a location in the same order as every other core.  On a
  shared bus, caches snoop each other's operations, and the bus orders them.
  Write-update coherence broadcasts every write to all copies;
  write-invalidate coherence broadcasts one invalidation per block and wins
  for the common access patterns.  MSI
  implements write-invalidate with three states; cache-to-cache transfers
  serve reads from the cache that holds the current value, the O state of
  MOSI avoids writing that value to memory at every transfer, and the E
  state of MESI avoids invalidations for data used by a single core; MOESI
  combines both.  Machines without a shared bus replace broadcasts by a
  directory that records which caches hold each block and orders the
  requests for it.  Invalidation causes coherence misses, the fourth kind of
  miss, which stem from true or false sharing.
\end{keypoint}

\end{document}
